\section{Requerimientos del problema}

Los requerimientos funcionales indican lo que debe hacer el sistema que resuelva el problema y los no funcionales, las cualidades que debe tener. Ninguno fija un algoritmo; la columna de origen remite al caso de uso, dato o restricción del que se derivan, y la Tabla~\ref{tab:trazabilidad-rf} completa su relación con actores, datos y criterios. Se marca como alta la prioridad de los que, de incumplirse, invalidarían las conclusiones: la etiqueta, la fuga de información y la evaluación temporal.

\begin{longtable}{p{1.2cm}p{7.0cm}p{2.4cm}p{1.6cm}p{1.2cm}}
\caption{Requerimientos funcionales y no funcionales}
\label{tab:requerimientos}\\
\toprule
Id & Requerimiento & Origen & Obj. & Prior. \\
\midrule
\endfirsthead
\multicolumn{5}{l}{\footnotesize\itshape (continuación)}\\
\toprule
Id & Requerimiento & Origen & Obj. & Prior. \\
\midrule
\endhead
\bottomrule
\endfoot
RF-01 & El sistema debe contrastar los totales de siniestros, lesionados y fallecidos del archivo con los reportes oficiales de un periodo común & CU7; M1, M5 & General & Alta \\
RF-02 & El sistema debe construir la etiqueta con la definición y la fecha de corte de la sección~2.2 & CU6, CU8; \texttt{FALLECIDOS} & General & Alta \\
RF-03 & El sistema debe excluir de las entradas toda variable derivada de las víctimas y registrar la lista final de variables & CU8; M2, M4 & General, OE2 & Alta \\
RF-04 & El sistema debe tratar los valores faltantes, \texttt{ND} y las categorías inconsistentes con reglas explícitas y documentadas & CU8; sección~6.3 & OE1, OE2 & Media \\
RF-05 & El sistema debe reportar por grupo la proporción de graves, el número de casos y una medida de incertidumbre & CU5; conjunto A & OE1 & Alta \\
RF-06 & El sistema debe evaluar la clasificación con casos posteriores en el tiempo a los usados para obtenerla & CU6; M4 & OE2 & Alta \\
RF-07 & El sistema debe comparar al menos cinco tipos de modelos de aprendizaje automático clásico con dos referencias sencillas, con los conjuntos A y B & CU6, CU9; R-03 & OE2 & Alta \\
RF-08 & El sistema debe medir su desempeño sobre la clase grave y no con la exactitud global, que alcanza el 90,4\,\% sin identificar ningún siniestro grave & CU6; R-06 & OE2 & Alta \\
RF-09 & El sistema debe indicar si las probabilidades que estima corresponden a las proporciones observadas & CU6; etiqueta & OE2 & Media \\
RF-10 & El sistema debe presentar las circunstancias más asociadas con la gravedad de forma comprensible para usuarios no especialistas & CU5, CU6 & OE1, OE2 & Media \\
RF-11 & El sistema debe permitir consultar un siniestro registrado y devolver la probabilidad estimada, las circunstancias asociadas y una advertencia de uso retrospectivo & CU6; modelo seleccionado & OE3 & Media \\
RNF-01 & El sistema debe producir las mismas cifras a partir del mismo archivo y la misma versión del procedimiento (reproducibilidad) & CU7; Anexo~B & General, OE2 & Alta \\
RNF-02 & El sistema debe evaluarse con casos no usados para obtener la estimación y con medidas apropiadas al desbalance (confiabilidad) & CU6; etiqueta & OE2 & Alta \\
RNF-03 & El sistema debe registrar toda limpieza o exclusión de datos (trazabilidad) & CU8; sección~6 & OE1 & Media \\
RNF-04 & El sistema debe permitir explicar cada circunstancia que señale como asociada a la gravedad (interpretabilidad) & CU5, CU6 & OE2, OE3 & Media \\
RNF-05 & El sistema debe publicar solo resultados agregados y nunca registros individuales (privacidad) & CU5, CU6; R-07 & OE1 & Alta \\
\end{longtable}
